\chapter{Data and Processing Pipeline}

\section{Market and news data}

The analytical market file contains 67,035 observations for 85 series: 83 company shares plus PX1 and TBIDX. Stored
dates run from 17 June 2021 to 12 October 2026. Each row contains a company
name, symbol, date, opening price, high, low, closing price, and volume. The
per-symbol files remain protected inputs; the combined full-universe file is a
derived analytical input.

The combined news file contains 12,498 deduplicated records for 90 symbols,
with valid dates from 29 February 2012 to 12 October 2026. Each row contains a
symbol, date, time, headline, and source URL. The news and market universes are
not identical because their collection processes include different sets of
symbols.

\begin{table}[htbp]
\centering
\caption{Stored analysis datasets}
\begin{tabularx}{0.94\textwidth}{@{}lrrX@{}}
\toprule
Dataset & Rows & Symbols & Fields and date range \\
\midrule
Market & 67,035 & 85 & Name, symbol, date, OHLCV; 17 Jun 2021--12 Oct 2026 \\
News & 12,498 & 90 & Symbol, date, time, headline, URL; 29 Feb 2012--12 Oct 2026 \\
\bottomrule
\end{tabularx}
\end{table}

Symbols provide the join key between the two sources. A market row can exist
without a nearby headline, and a news symbol can exist without a usable market
series. The classification stage therefore permits unmatched records and does
not substitute one symbol for another.

\section{Pipeline structure}

The default runner performs analysis on local files. It does not scrape and it
does not call an external generative service. The principal stages are shown
in Figure~\ref{fig:long-pipeline}.

\begin{figure}[htbp]
\centering
\begin{tikzpicture}[
  node distance=6mm and 9mm,
  box/.style={draw=blue,rounded corners,fill=softblue,align=center,
    minimum width=33mm,minimum height=10mm,font=\small},
  arrow/.style={-{Latex[length=2.4mm]},thick,color=blue}
]
\node[box] (raw) {Protected ticker\\and news files};
\node[box,right=of raw] (combine) {Validated 85-series\\analysis file};
\node[box,right=of combine] (detect) {Spike and decay\\detectors};
\node[box,below=of detect] (classify) {News matching\\and categories};
\node[box,left=of classify] (watch) {Watchlist and\\market refinement};
\node[box,left=of watch] (eval) {Evaluation, manifest,\\and summaries};
\node[box,below=of watch] (serve) {Flask APIs and\\React Terminal};
\draw[arrow] (raw) -- (combine);
\draw[arrow] (combine) -- (detect);
\draw[arrow] (detect) -- (classify);
\draw[arrow] (classify) -- (watch);
\draw[arrow] (watch) -- (eval);
\draw[arrow] (watch) -- (serve);
\draw[arrow] (eval) |- (serve);
\end{tikzpicture}
\caption{File-based processing and presentation pipeline.}
\label{fig:long-pipeline}
\end{figure}

The file-based design makes intermediate results inspectable. An analyst can
compare a watchlist row with its classified flag and then with the originating
market values. The evaluation manifest records normalized text hashes for the
detector, evaluation configuration, event register, analytical inputs, and
generated evaluation files.

\section{Preprocessing and data quality}

Dates are parsed and observations are sorted by symbol and date. TBIDX, an
index series, and PLTU, a known incomplete ticker series, are excluded from the
spike detector. The analytical input still retains these symbols for other
uses. Percentage return is calculated from consecutive stored closes.

A return is set to missing when the gap between consecutive observations is
longer than ten calendar days. This rule prevents a multi-week or multi-month
change from being treated as one daily movement. The return baseline then uses
the most recent 60 valid prior returns, so one omitted gap return does not
remove scores from the following 60 rows.

Volume is transformed with $\log(1+V)$, which keeps zero volume defined and
reduces the influence of very large raw values. Source open, high, low, close,
and volume fields remain unchanged. Rows with invalid ranges, missing or
negative values, or prices outside the recorded range are written to a
separate quality file for review.

\section{Derived outputs}

The spike and decay detectors write separate outputs. Classification adds
headline, timing, and category fields, while the watchlist joins unexplained
combined-signal rows with decay cases that lack a news match. Refinement adds
index movement, same-day co-flags, excess return, and a market-context tag.

The evaluation directory stores event results, metrics, baselines, sensitivity
and funnel counts, figures, a summary, and a manifest. The Terminal reads but
does not rewrite these files. Each output can be regenerated, while protected
inputs remain unchanged; dashboard values present the stored analysis rather
than a second calculation.
